% Paper B (problem131_multistate.tex): material removed from the paper on
% 2026-10-02 during the second arXiv edit (journal-style and length pass).
% Each piece is verbatim from the file and lines named in its header, as they
% were before this edit. Labels refer to that version, so this file is not
% meant to be compiled. No theorem, proposition, lemma or corollary lost
% its proof (the last piece is a second proof of a lemma that keeps its first).

%----------------------------------------------------------------------
% Source: paperB/simplex_thresholds.tex, lines 154-159.
% Sentence on the inverse Hessian of the rate (Paper C, in preparation) after Theorem thm:simplex-interior. Replaced by one sentence.
%----------------------------------------------------------------------
$\alpha$. In a paper in preparation~\cite{paperC} we show that the inverse
Hessian of the rate is the occupation covariance of a tilted chain, on every
face where the chain is irreducible. So $D(\alpha)$ should be a Gaussian
factor. We now show this directly and write it
as a spanning-tree sum. The next lemma computes the determinant of an occupation
covariance for any reversible chain.

%----------------------------------------------------------------------
% Source: paperB/simplex_thresholds.tex, lines 240-247.
% Reading of the factors of Proposition prop:kirchhoff(3) for the uniform start. Cut, no proof uses it.
%----------------------------------------------------------------------
The proof is a direct check, given in Appendix~\ref{app:thresholds}. Here is how
to read the factors in (3) for the uniform start. First,
$A^2=d(\mu\cdot h)(h\cdot\mathbf1)$, one factor for the start of the path and one
for its free end, and Proposition~\ref{prop:start-weights} explains the start factor.
Next, $(4\pi)^{-(k-1)/2}$ is the Gaussian normalization $(2\pi)^{-(k-1)/2}$ times
$2^{-(k-1)/2}$, from the $2$ in the variance $2\langle f,(-G)^{-1}f\rangle_\pi$.
And $\sqrt{\tree(\alpha)}/\prod_i\alpha_i=2^{(k-1)/2}(\det\Sigma_\alpha)^{-1/2}$ by
Lemma~\ref{lem:occupation-covariance}.

%----------------------------------------------------------------------
% Source: paperB/general_start.tex, lines 164-175.
% Heuristic next-term prediction of the balancing gain with a start (not proved). Cut, no proof uses it.
%----------------------------------------------------------------------
This example only shows that balancing can fail.
Theorem~\ref{thm:start-maxima} below shows that a start gains at most a factor
$1+O(1/L)$ over the balanced bins. For every $\mu$ that is not constant on $F$, a
heuristic next-term expansion predicts more. For all large $N$ and
$cL\le T\le CL$, it predicts that the largest bin on $F$ beats the balanced bins
by a factor
$1+\frac1{4\mu(F)^2T}\sum_{a\in F}(\mu_a-\mu(F)/k)^2+o(1/T)$. So balancing
should fail there. The formula comes from expanding $\log S_n(u)$
(Theorem~\ref{thm:simplex-interior}) and the start weights $\sqrt{\alpha_a}/A$
to second order around the center of $F$ and maximizing over $n$. We do not
prove this, since the errors of the start weights are not controlled at that
order.

%----------------------------------------------------------------------
% Source: paperB/general_start.tex, lines 242-259.
% Remark rem:start-lines in full. The criterion at t^* (an iff with its proof) was cut. The monotonicity and the definition of t^* stay.
%----------------------------------------------------------------------
\begin{remark}[the order of the winners]\label{rem:start-lines}
The winning size is nondecreasing in $t$. Indeed, if $k_1$ attains
$\max_k\ell^\mu_k$ at $t_1$ and $k_2$ attains it at $t_2>t_1$, adding the two
inequalities $\ell^\mu_{k_1}(t_1)\ge\ell^\mu_{k_2}(t_1)$ and
$\ell^\mu_{k_2}(t_2)\ge\ell^\mu_{k_1}(t_2)$ gives $(k_2-k_1)(t_2-t_1)\ge0$. Put
$t^*=a_d+\frac1{d-1}\log(d\mu_{(1)})$, where $\ell^\mu_1$ and $\ell^\mu_d$ cross.
Some size $2\le k\le d-1$ is the unique winner on an open interval of $t$ if and
only if $\ell^\mu_k(t^*)>\ell^\mu_1(t^*)$ for some such $k$. Indeed, suppose
that $\ell^\mu_k(t^*)\le\ell^\mu_1(t^*)=\ell^\mu_d(t^*)$ for every such $k$.
The slope $k-1$ of $\ell^\mu_k$ lies strictly between the slopes $0$ and $d-1$,
so $\ell^\mu_k<\ell^\mu_1$ for $t<t^*$ and $\ell^\mu_k<\ell^\mu_d$ for $t>t^*$.
Conversely, if some intermediate line is above $\ell^\mu_1$ at $t^*$, then among
the lines that are highest at $t^*$ the one with the largest slope is
intermediate, and by continuity it is the unique winner on an interval
$(t^*,t^*+\eta)$. For the uniform start $t^*=a_d$ and
$\ell^\mu_k(t^*)-\ell^\mu_1(t^*)=(k-1)(a_d-a_k)<0$, which agrees with
Corollary~\ref{cor:simplex-modes}.
\end{remark}

%----------------------------------------------------------------------
% Source: paperB/general_start.tex, lines 297-320.
% Remark rem:K3-finite in full, before it was shortened (the two heuristic window ends as displays, the gain 1/(24z), the width 0.0260580-0.0546957/L).
%----------------------------------------------------------------------
\begin{remark}[the window at finite $N$]\label{rem:K3-finite}
The balanced bins of the edge cross the vertices exactly at Paper~A's root
$u_N$, for every $N\ge2$, since $\mu$ is constant on $\{1,2\}$. Indeed
$P^\mu_N(n_1,n_2,0)/P^\mu_N(N,0,0)=S_{(n_1,n_2)}(u)$, which is the polynomial of
Paper~A~\cite{paperA} at $u=r/p$. This crossing is at $T=Nu_N/(1+2u_N)$, which
differs from Paper~A's $z_N=Nu_N/(1+u_N)$ by $O(L^2/N)$. Near $t=a_2$ the lines
of the full face and of the edges through $3$ stay below those of the vertices,
so for large $N$ this crossing is the lower end of the window. For the 2 ends at
finite $N$ we use $z=Nu=T/p$. A heuristic calculation, which we do not prove,
gives
\[
 z-L+\tfrac12\log z=a_2+\frac1{8z}+\frac{z^2}{2N}+\cdots,\qquad
 z-L+\tfrac12\log z=t_{\rm ef}+\frac1{12z}+\frac{3z^2}{2N}+\cdots,
\]
where only $a_2+\frac1{8z}$ is proved, by Theorem~\ref{thm:simplex-crossing}.
The terms $z^2/N$ come from a discrete term $-\frac{(k-1)^2}2\frac{z^2}N$ in the
logarithm of~\eqref{eq:simplex-uniform-ratio}, and the upper end also uses a gain
$\frac1{24z}$ of the best full-face bin over the balanced one. Without the terms
$z^2/N$ the width is $0.0260580-0.0546957/L+\cdots$. Exact computations give
the widths $0.0843$, $0.0324$, $0.0223$ and $0.0214$ at $N=300$, $3000$,
$30000$ and $300000$, and the prediction gives $0.0802$, $0.0321$, $0.0226$ and
$0.0217$ (numerically verified, see Section~\ref{sec:verify}). Below $N$ of
about $10^4$ the terms $z^2/N$ are larger than the terms $1/L$.
\end{remark}

%----------------------------------------------------------------------
% Source: paperB/frontier_modes.tex, lines 114-144.
% Remark rem:frontier-superlevel (proof omitted) and the paragraph on quasi-ergodic laws after it. Cut, no proof uses them.
%----------------------------------------------------------------------
\begin{remark}[how many bins beat a vertex]\label{rem:frontier-superlevel}
The same profile counts the bins that beat a vertex. Put $h_*=\max_ih_i$ and,
for a support $S$ with $|S|=k\ge2$, $B_S(t,h)=\ell_S(t,h)-h_*$. Let $M_{N,S}(t,h)$
be the number of bins with support exactly $S$ whose tilted probability is
strictly greater than that of every vertex. Then, uniformly for $(t,h)$ in
compact sets,
\[
 \frac{L^{(k-1)/2}}{N^{k-1}}M_{N,S}(t,h)
 \longrightarrow
 \frac{|\mathbb B_{k-1}|}{\sqrt k}
 \left(\frac{4\,[B_S(t,h)]_+}{k^2}\right)^{(k-1)/2},
\]
where $|\mathbb B_{k-1}|$ is the volume of the unit ball in $\mathbb R^{k-1}$
and $[x]_+=\max(x,0)$ (proof omitted). In the coordinates
$y=\sqrt L\,(n/N-\frac1k\mathbf1)$ these bins fill a ball of squared radius
$4B_S/k^2$, and $\sqrt k\,(\sqrt L/N)^{k-1}$ is the covolume of the lattice of
the points $y$. For $k=2$ the constant is $\sqrt{2[t-a_2]_+}$, which agrees with
the binary window of Paper~A once the different centering is taken into
account. The profile and the lattice volumes are checked numerically (see
Section~\ref{sec:verify}), and Lean checks that the Gram determinant of the
vectors $e_i-e_k$, $i<k$, is $k$.
\end{remark}

On $K_d$ the quasi-ergodic law of a face, the product of its left and right
Perron vectors, is uniform on the face, and $\bar h_S$ is the average
of $h$ under it. For a general generator $Q$ we have
$P_N\operatorname{diag}(e^{h/N})=I+\frac TN(Q+\operatorname{diag}(h)/T)+O(T/N^2)$,
so the field acts as a potential $h/T$ added to $Q$. First-order perturbation of
the Perron root then lowers $T\lambda_F$ by $\langle h,\alpha^*_F\rangle+O(1/T)$,
where $\alpha^*_F$ is the quasi-ergodic law. So we expect each face line to gain
$\langle h,\alpha^*_F\rangle$ in general, but we have not written this out.

%----------------------------------------------------------------------
% Source: paperB/frontier_modes.tex, lines 157-170.
% Remark rem:start-vs-field in full (a shortened version was later cut as well).
%----------------------------------------------------------------------
\begin{remark}[start versus field]\label{rem:start-vs-field}
Compare a start $\mu$ with full support and no field
(Corollary~\ref{cor:start-modes}) with the uniform start and the field
$h_i=\log(d\mu_i)$ (Theorem~\ref{thm:weakfield-phases}). We add $\log d$ to the
start lines, which does not change the modes. Then a vertex $i$ scores
$\log(d\mu_i)$ in both. A support $S$ with $k\ge2$ states scores $(k-1)(t-a_k)$
plus the logarithm of the arithmetic mean of $d\mu$ on $S$ under the start, and
plus the logarithm of the geometric mean under the field. The arithmetic mean is
at least the geometric mean, with equality only when $\mu$ is constant on $S$.
So a start favors mixed faces more than the field with the same vertex scores.
For example, in Example~\ref{ex:K3-start} state $3$ has start weight $0$, which
as a field would mean $H=\infty$, and yet the edge window has width only
$0.026$. We have not written out the lines for a start and a field together.
\end{remark}

%----------------------------------------------------------------------
% Source: paperB/frontier_modes.tex, lines 202-210.
% Remark rem:mass-mode. Cut, no proof uses it.
%----------------------------------------------------------------------
\begin{remark}[mass versus mode]\label{rem:mass-mode}
From each state of a face $F$ with $k$ states the chain stays in $F$ with
probability $1-(d-k)r$, so $\Prob(\supp K_N\subseteq F)=\mu(F)(1-(d-k)r)^{N-1}$.
Every bin without full support lies on a face with $d-1$ states. So the bins
with full support have mass at least
$1-\sum_{|F|=d-1}\mu(F)(1-r)^{N-1}=1-O(e^{-T})$, while for the uniform start,
$t<a_d$ and large $N$ the modes are single states
(Corollary~\ref{cor:simplex-modes}).
\end{remark}

%----------------------------------------------------------------------
% Source: paperB/frontier_boundary.tex, lines 155-161.
% Last paragraph of Subsection sec:vertex-completion (results with proofs omitted, checked numerically). Cut, no proof uses it.
%----------------------------------------------------------------------
We have also found the next term of $u_{b,a}$ for fixed rare counts, which
for one rare count is the 2-state fixed-edge expansion of Paper~A, the growth
$u_{b,a}\sim\ell\log|a|/(2A\sqrt{|a|b})$ with $A=\sum_i\sqrt{a_i/|a|}$ when
$|a|\to\infty$ with proportions $a_i/|a|$ bounded away from $0$, and an example
with $|a|=o(b)$ where $S_{(b,a)}(u)\sim H(bu^2)$ fails although the root is
still correct (proofs omitted, checked numerically, see
Section~\ref{sec:verify}).

%----------------------------------------------------------------------
% Source: paperB/app_boundary.tex, lines 285-348.
% The remark after Theorem thm:simplex-boundary-crossover in full. It was shortened to Remark rem:crossover-numbers, and the description of verify_crossover_factor.py moved to the verification section.
%----------------------------------------------------------------------
\begin{remark}
The 3 terms of $\mathcal E$ have different sources. Replacing $M$ by $N$ in
the exponent $A^2Y$ gives a factor $\exp(-A^2z\sum_ja_j/N)$. The Gaussian
average of the factors $\Phi$ gives back $\exp(z\sum_ja_j/N)$, which is the
term $v\sum_ja_j$ in \eqref{eq:simplex-boundary-M-form}. Together they give
the first term of $\mathcal E$. The second term comes from the same
Gaussian average. The factors $\Phi$ and the power of $w$ shift the saddle
point, and $\psi$ measures this shift. The last term does not depend on the
sizes $a_j$. For $\ell=0$ it is all of $\mathcal E$, and for equal $\beta_i$
its coefficient $\gamma_0$ is then $-(k-1)/8$, the coefficient of $1/z$ in
\eqref{eq:simplex-uniform-ratio}. For $\ell=1$ the first term
of $\mathcal E$ dominates for large $x_1$. For $\ell\ge2$ the cross terms
are also linear in the $x_j$, since $\psi(x)/\sqrt x=I_2(2\sqrt x)/I_1(2\sqrt x)$
tends to $1$ by \cite[Eq.~10.40.1]{dlmf}.

Take 2 equal large counts and one rare count, so that $g=1$ and $\eta=5/2$,
and let $z=\log N$. Then $x_1=2\rho$, and the first 2 terms of
$z\mathcal E$ are $-\rho+1.25\,\psi(2\rho)$. They nearly cancel at
$\rho=1$, so the first term alone is not the correction. The ratio of
$S_{(n,a)}$ to the right side of \eqref{eq:simplex-boundary-profile}
without $e^{\mathcal E}$ is $1.04$ for $\rho=1$ and $N=2000$, and with the
first term alone it would be $1.19$. For $\rho=8$
it is $0.36$, $0.53$, $0.63$ and $0.70$ at $N=2\cdot10^3$, $2\cdot10^4$,
$2\cdot10^5$ and $2\cdot10^6$. With the factor $e^{\mathcal E}$ the value
for $\rho=1$ becomes $0.90$, and the 4 values for $\rho=8$ become $0.52$,
$0.70$, $0.79$ and $0.85$. At $N=2000$ we have $L=7.6$, and the terms of
order $L^{-2}$ and $L^2/N$ are comparable to $\gamma_0/L$, so the value for
$\rho=1$ does not improve. The form \eqref{eq:simplex-boundary-M-form} does
not replace $\xi_j$ by $x_j$ or $Y^\eta M^{-(k-1)}$ by $z^\eta N^{-(k-1)}$.
It gives $0.87$ for $\rho=1$, and $0.78$, $0.92$, $0.96$ and $0.97$ for
$\rho=8$.

The constant in $O(L^{-2})$ still depends on $R$, and the proof gives only
a crude bound for it. At
$N=10^{12}$, with $z=L$, the quantity $L^2$ times the logarithm of the ratio
for \eqref{eq:simplex-boundary-profile} is about $-1.05$, $-2.6$, $-14$,
$-36$ and $-266$ for $\rho=1/4$, $1$, $4$, $8$ and $32$. For
\eqref{eq:simplex-boundary-M-form} it is about $-0.9$, $-1.4$, $-3.5$,
$-6.2$ and $-21$. Expanding $(1-s_a)^{\eta-k+1}$ and $\Phi(x_j(1-s_a))$ to
first order in $s_a=\sum_j\rho_j/L^2$ shows that the difference of the 2
constants is
$-(\sum_j\rho_j)(\eta-k+1+\sum_j\psi(x_j))+O(L^{-1}+L^4/N)$. Here this
is $-\rho(\tfrac12+\psi(2\rho))$, about $-248$ for $\rho=32$. So the first
constant grows roughly like $\rho^{3/2}$ in these examples, while the
second grows roughly like $\rho$.

The expansion \eqref{eq:simplex-boundary-root} converges slowly. In the
same example, for $\rho=8$ and $N=10^6$, it gives $6.58$, while
$Nu_{n,a}=14.99$. The root $\zeta_N$ of
\eqref{eq:simplex-boundary-implicit-root} is $14.86$, and the root of
\eqref{eq:simplex-boundary-M-form} is $14.97$. For $\rho=1$ the 4 values
are $15.41$, $17.86$, $17.85$ and $17.86$. The loss is caused by expanding
around $z=cL$. Here $c=2$, while $Nu_{n,a}/L$ is only $1.08$ for $\rho=8$
and $1.29$ for $\rho=1$.

These numbers come from \texttt{verify\_crossover\_factor.py}. For
$N\le2\cdot10^6$ it computes $S_{(n,a)}$ in 3 ways: by the proper-run formula
and by a quadrature of \eqref{eq:boundary-integral-form}, both in double
precision, and by the sum \eqref{eq:simplex-positive-sum} in 60-digit decimal
arithmetic. The values at $N=10^{12}$ and the crossings come from the
quadrature and the sum only. These agree to about $10^{-14}$ in $\log S$, and
the crossings, found by bisection with the sum, agree with the quadrature to
about $10^{-7}$.
\end{remark}

%----------------------------------------------------------------------
% Source: paperB/general_start.tex, lines 139-142 (third edit, 2026-10-02).
% Definition of the face constant K_F after Theorem thm:start-crossings. Cut, K_F is not used later.
%----------------------------------------------------------------------
Call $K_F=\mu(F)k^{k-1/2}(4\pi)^{-(k-1)/2}$ the second-order constant of a face
$F$ with $k\ge2$ states, and put $K_{\{i\}}=\mu_i$ for a vertex. Then
$a_k(\mu)=-\frac1{k-1}\log(K_F/K_{\{i\}})$, so $K_F$ depends on the start only
through $\mu(F)$.

%----------------------------------------------------------------------
% Source: paperB/simplex_thresholds.tex, lines 191-195 (third edit, 2026-10-02).
% Second (Pruefer code) proof of Lemma lem:weighted-cayley. Cut, the first proof (matrix determinant lemma and matrix-tree theorem) stays.
%----------------------------------------------------------------------
For a
second proof, note that a tree $\mathcal T$ has weight
$\prod_ix_i^{\deg_{\mathcal T}(i)}$. The Pr\"ufer code is a bijection from the
labeled trees on $\{1,\ldots,k\}$ to the words of length $k-2$ in which $i$
occurs $\deg_{\mathcal T}(i)-1$ times, so the sum is $\prod_ix_i\cdot X^{k-2}$.
